\addstarredchapterc{Abstract in Greek}
\chapter*{Περίληψη}
\begin{abstractGr}
  Η ικανότητα παροχής αδιάλειπτης συνδεσιμότητας και χρήσης αρχιτεκτονικής ``cloud-first'' (με προτεραιότητα στο cloud) για τα Συστήματα Διαχείρισης Παραγγελιών και τον κλάδο των Εκδηλώσεων έχει γίνει κύριο σημείο εστίασης της ψηφιακής ανατροπής (digital disruption) και στους δύο αυτούς τομείς. Ωστόσο, λόγω της φύσης των προσωρινών τους εκδηλώσεων (π.χ. Πολιτιστικά Φεστιβάλ) και των σχετικών ad-hoc υποδομών δικτύου, τα περιβάλλοντα αυτά θεωρούνται εξαιρετικά εχθρικά. Αυτό οφείλεται σε μεγάλο βαθμό σε ακραίες εξάρσεις καθυστέρησης (latency spikes), σε τοπικές Παρεμβολές Ραδιοσυχνοτήτων (RFI) και στην πλήρη απώλεια συνδεσιμότητας. Ως εκ τούτου, σε περιβάλλοντα όπου υπάρχουν πολυάριθμα σημεία υψηλού φόρτου (π.χ. τερματικά POS, Κιόσκια, κ.λπ.), η χρήση της Σύγχρονης Περιοδικής Ερώτησης HTTP (Synchronous HTTP Polling) στα Παραδοσιακά Συστήματα Σημείων Πώλησης (POS) συχνά προκαλεί προβλήματα που αναφέρονται ως το ``Πρόβλημα των Δύο Στρατηγών'', το οποίο οδηγεί σε ακραίο Αποσυγχρονισμό Δεδομένων, Παραγγελίες-Φαντάσματα και Οικονομικές Ασυμφωνίες.
  
  Αυτή η πτυχιακή εργασία περιγράφει τον σχεδιασμό, την ανάπτυξη και την πειραματική επικύρωση ενός Συστήματος Διαχείρισης Παραγγελιών (OMS) που παρέχει ικανότητα υψηλής απόδοσης για τη διαχείριση παραγγελιών σε ένα τοπικό περιβάλλον. Το OMS είναι ειδικά σχεδιασμένο ώστε να είναι ικανό να διαχειρίζεται παραγγελίες σε εχθρικά περιβάλλοντα δικτύωσης. Ένα σημαντικό αρχιτεκτονικό στοιχείο αυτού του συστήματος χρησιμοποιεί μια Βελτιστοποιημένη, Καθοδηγούμενη από Γεγονότα (Event-Driven) Εσωτερική Τοπολογία, κατασκευασμένη πάνω στη στοίβα τεχνολογιών MERN (MongoDB, Express.js, React, Node.js). Ένα σημαντικό πλεονέκτημα της χρήσης αυτού του τύπου εσωτερικής τοπολογίας είναι ότι επιτρέπει την αποσύνδεση της Επιχειρησιακής Διεπαφής από τις Εξωτερικές Εξαρτήσεις του Cloud. Για να βελτιώσει περαιτέρω τις δυνατότητες απόδοσής της σε εχθρικά περιβάλλοντα δικτύωσης, η Προτεινόμενη Αρχιτεκτονική αντικαθιστά το Synchronous HTTP Polling με Αμφίδρομες (Full-Duplex) Συνδέσεις WebSocket και χρησιμοποιεί Ανθεκτικές Διεπαφές Virtual DOM βασισμένες σε React καθώς και Προσωρινή Αποθήκευση Offline-First (με προτεραιότητα στην εκτός σύνδεσης λειτουργία). Αυτό γίνεται για να εξασφαλιστεί ο Ντετερμινιστικός Συγχρονισμός Κατάστασης σε Απομονωμένους Επιχειρησιακούς Κόμβους (Ταμίες, Προσωπικό Κουζίνας, Διοικητικά Κέντρα). Επιπλέον, χρησιμοποιούνται και επιβάλλονται αυστηρά όρια ασφαλείας με τη χρήση JWTs και RBAC Μηδενικής Εμπιστοσύνης (Zero Trust) για τον Έλεγχο Πρόσβασης.

  Προκειμένου να αξιολογηθούν εμπειρικά τα χαρακτηριστικά απόδοσης της προτεινόμενης αρχιτεκτονικής, πραγματοποιήθηκαν αυστηρές αξιολογήσεις. Αυτές οι αξιολογήσεις περιλάμβαναν Δοκιμές Ακραίας Καταπόνησης (Stress Testing) Υψηλής Έντασης (Επικυρωμένη Σταθερότητα σε 1,58 Εκατομμύρια Ταυτόχρονες Συναλλαγές) και Προσομοιώσεις Μηχανικής Χάους (Chaos Engineering) που κατέδειξαν Στόχο Χρόνου Ανάκαμψης (RTO) από Σοβαρή Βλάβη στα 310 Χιλιοστά του δευτερολέπτου. Οι εργαστηριακές μετρήσεις που έγιναν κατά τη διάρκεια των αξιολογήσεων επικυρώθηκαν στη συνέχεια μέσω εκτεταμένων εφαρμογών σε πραγματικές συνθήκες (field deployments) κατά τη διάρκεια επτά ημερών με μαζικές ριπές κίνησης δεδομένων (Burst Traffic). Η τηλεμετρία που συλλέχθηκε κατά τη διάρκεια της εφαρμογής σε πραγματικές συνθήκες επικύρωσε ότι οι Ασύγχρονες Αρχιτεκτονικές Βρόχου Γεγονότων (Event Loop) μειώνουν σημαντικά τον φόρτο δεδομένων δικτύου (Bandwidth Payload) και εξαλείφουν την Υποβάθμιση Βάσει Σύνδεσης (Connection-Based Degradation) όταν εφαρμόζονται σε Εχθρικά Περιβάλλοντα Δικτύωσης.
  
  \vspace{1em}
  \textbf{Λέξεις Κλειδιά:} Σύστημα Διαχείρισης Παραγγελιών, Αρχιτεκτονική Καθοδηγούμενη από Γεγονότα, WebSockets, Μοντέλο Εκτέλεσης σε Τοπικό Επίπεδο (Local-First Topology), Δίκτυα Ad-hoc, Επεξεργασία Υψηλού Ρυθμού, Συγχρονισμός Πραγματικού Χρόνου.
\end{abstractGr}
